\section{Yêu cầu hệ thống (Requirements)}

\subsection{Yêu cầu chức năng (Functional requirements)}

\subsubsection{FR-1: Thu thập dữ liệu (Data ingestion)}

\begin{itemize}
  \item \textbf{FR-1.1} Hệ thống phải thu thập được dữ liệu từ mã nguồn (source code), tài liệu (SRS, v.v), và git log.
  \item \textbf{FR-1.2} Hệ thống phải đọc và ingest được các định dạng tài liệu phổ biến: Markdown và PDF.
  \item \textbf{FR-1.3} Hệ thống phải nạp dữ liệu từ các source được cấu hình, mỗi source thuộc một trong hai loại: (a) một Git repository (qua remote URL hoặc đường dẫn tới một repo trên đĩa) tại một ref cấu hình được, mặc định là default branch; và (b) một thư mục filesystem không yêu cầu version control. Source loại (a) cung cấp thêm \texttt{ref} và \texttt{commit\_sha} cho provenance (FR-1.4); loại (b) định danh phiên bản tài liệu qua \texttt{content\_hash}.
  \item \textbf{FR-1.4} Hệ thống phải gắn metadata nguồn gốc cho mỗi đơn vị dữ liệu được nạp. Provenance bắt buộc gồm \texttt{source\_id}, đường dẫn file, \texttt{content\_hash} và thời điểm index (\texttt{indexed\_at}); riêng nguồn Git bổ sung \texttt{ref} (branch) và \texttt{commit\_sha}.
  \item \textbf{FR-1.5} Hệ thống phải cho phép cấu hình danh sách nguồn cần nạp, bao gồm việc loại trừ (ignore) các file hoặc thư mục không cần index.
\end{itemize}

\subsubsection{FR-2: Xử lý \& lập chỉ mục (Processing \& indexing)}

\begin{itemize}
  \item \textbf{FR-2.1} Hệ thống phải chia dữ liệu thành các chunk trước khi sinh embedding, cắt theo ranh giới cấu trúc tự nhiên và không cắt đôi một đơn vị không thể chia: mã nguồn theo ranh giới khai báo (hàm, lớp), tài liệu theo cấu trúc văn bản (heading, đoạn).
  \item \textbf{FR-2.2} Hệ thống phải sinh embedding cho từng chunk thông qua embedding provider đang được cấu hình.
  \item \textbf{FR-2.3} Hệ thống phải xây dựng Knowledge Graph gồm các entity và quan hệ giữa chúng (theo FR-3). Entity gồm: các entity mã nguồn theo các mức ở FR-3.4, cùng các entity phi mã nguồn là tài liệu, yêu cầu (SRS) và commit.
  \item \textbf{FR-2.4} Hệ thống phải lưu trữ đồng thời vector embedding, Knowledge Graph và metadata, sao cho mỗi chunk truy ngược được về nguồn gốc.
\end{itemize}

\subsubsection{FR-3: Liên kết tri thức (Knowledge linking)}

\begin{itemize}
  \item \textbf{FR-3.1} Hệ thống phải tạo được các quan hệ cấu trúc và phụ thuộc trong mã nguồn, bao gồm: chứa/khai báo (CONTAINS/DECLARES), gọi hàm (CALLS), nhập dependency (IMPORTS), kế thừa (EXTENDS), hiện thực interface (IMPLEMENTS) và ghi đè (OVERRIDES).
  \item \textbf{FR-3.2} Hệ thống có thể bổ sung các quan hệ sâu hơn: khởi tạo (INSTANTIATES), đọc/ghi biến, field hoặc constant (READS/WRITES), tham chiếu một định danh khác mà không gọi hàm (REFERENCES), dùng kiểu (HAS\_TYPE), annotation (ANNOTATED\_WITH), ném ngoại lệ (THROWS), liên kết test (TESTS).
  \item \textbf{FR-3.3} Hệ thống phải liên kết giữa các artifact khác nhau, gồm: tài liệu mô tả code (DESCRIBES), yêu cầu được hiện thực bởi code (IMPLEMENTED\_BY), yêu cầu được kiểm bởi test (VERIFIED\_BY), và commit sửa đổi code (MODIFIES). Việc liên kết dựa trên các tín hiệu: tên định danh (tên hàm, tên lớp), từ khóa, và tham chiếu đường dẫn; mỗi liên kết suy luận theo heuristic phải kèm một điểm tin cậy (confidence score) để phân biệt liên kết chắc chắn với liên kết phỏng đoán.
  \item \textbf{FR-3.4} Hệ thống phải biểu diễn tri thức theo các mức phân cấp, từ mịn đến thô: constant, field, function/method, class/interface, file, module, package, project; và quan hệ chứa giữa các mức (CONTAINS, theo FR-3.1).
  \item \textbf{FR-3.5} Hệ thống phải hỗ trợ các quan hệ nối entity thuộc các source khác nhau trong cùng một product (cross-source), áp dụng cho cả quan hệ cấu trúc (IMPORTS tới thư viện dùng chung) và quan hệ cross-artifact ở FR-3.3 (DESCRIBES, IMPLEMENTED\_BY, VERIFIED\_BY nối tài liệu/yêu cầu ở một source tới code ở source khác).
\end{itemize}

\subsubsection{FR-4: Truy xuất (Retrieval)}

\begin{itemize}
  \item \textbf{FR-4.1} Hệ thống phải nhận truy vấn ở dạng văn bản ngôn ngữ tự nhiên (free-text) và dùng trực tiếp làm đầu vào cho truy xuất (FR-4.2). Hệ thống không tự hiểu hay diễn giải ý định ngôn ngữ, đó là vai trò của agent.
  \item \textbf{FR-4.2} Hệ thống phải thực hiện hybrid search, kết hợp semantic search, keyword search và graph traversal, rồi hợp nhất (fusion) thành một danh sách kết quả.
  \item \textbf{FR-4.3} Hệ thống phải trả kết quả dưới dạng JSON có cấu trúc, mỗi kết quả kèm relevance score và metadata (nguồn, đường dẫn file, vị trí dòng code, ref/phiên bản).
  \item \textbf{FR-4.4} Hệ thống phải hỗ trợ các tham số truy vấn tùy chọn, bao gồm: giới hạn số kết quả (top-k) và lọc theo source, ref hoặc loại dữ liệu. Bộ lọc source chỉ được thu hẹp trong phạm vi các source mà principal được đọc, không bao giờ mở rộng.
  \item \textbf{FR-4.5} Khi truy vấn tham chiếu tới một ref chưa được index, hệ thống phải trả kết quả theo ref canonical và chỉ rõ điều này trong metadata.
\end{itemize}

\subsubsection{FR-5: API \& MCP}

\begin{itemize}
  \item \textbf{FR-5.1} Hệ thống phải cung cấp REST API cho cả thao tác truy vấn và thao tác quản trị (nạp, cập nhật, xóa nguồn và kiểm tra trạng thái).
  \item \textbf{FR-5.2} Hệ thống phải cung cấp một MCP server cung cấp các tool: một tool truy vấn tri thức và một tool kích hoạt cập nhật/đồng bộ, để AI Agent vừa tích hợp truy vấn vừa chủ động cập nhật trực tiếp.
  \item \textbf{FR-5.3} Hệ thống phải cung cấp tài liệu API mô tả request, response và ví dụ, đủ để agent hiểu cách dùng và diễn giải kết quả.
  \item \textbf{FR-5.4} Mọi lỗi API phải trả về theo một schema thống nhất gồm mã lỗi và thông điệp.
\end{itemize}

\subsubsection{FR-6: Đồng bộ \& cập nhật (Sync \& update)}

\begin{itemize}
  \item \textbf{FR-6.1} Việc cập nhật tri thức là trigger-based: hệ thống phải cung cấp thao tác cập nhật/đồng bộ qua REST API và MCP tool để agent hoặc caller chủ động kích hoạt khi nguồn thay đổi. Thao tác này phải idempotent (kích hoạt lại khi nội dung không đổi là một no-op).
  \item \textbf{FR-6.2} Khi được kích hoạt, hệ thống phải tự động nhận diện file mới, sửa đổi hoặc xóa và chỉ xử lý phần thay đổi (incremental), thay vì re-index toàn bộ.
  \item \textbf{FR-6.3} Hệ thống phải tránh xử lý trùng lặp: chunk có nội dung không đổi thì không được sinh lại embedding khi đồng bộ.
  \item \textbf{FR-6.4} Hệ thống phải gỡ bỏ (evict) tri thức lỗi thời khi nguồn tương ứng bị xóa hoặc thay đổi, giữ chỉ mục nhất quán với nguồn.
  \item \textbf{FR-6.5} Hệ thống phải cho biết độ mới của tri thức (qua \texttt{indexed\_at} và \texttt{commit\_sha}) trong kết quả truy vấn hoặc endpoint trạng thái, để agent quyết định có cần kích hoạt cập nhật trước khi truy vấn.
\end{itemize}

\subsubsection{FR-7: Cấu hình (Configuration)}

\begin{itemize}
  \item \textbf{FR-7.1} Hệ thống phải cho phép chọn embedding provider qua cấu hình (gọi API bên ngoài hoặc dùng local model).
  \item \textbf{FR-7.2} Hệ thống phải hỗ trợ cấu hình theo nhiều môi trường (dev, server, cloud).
\end{itemize}

\subsubsection{FR-8: Truy cập đồng thời \& bảo mật (Concurrency \& access)}

Các yêu cầu FR-8.2 - 8.13 dựa trên mô hình hai bước:
\begin{itemize}
  \item \textbf{authentication} resolve credential của request về một principal (\texttt{credential} $\rightarrow$ \texttt{principal\_id});
  \item \textbf{authorization} tính tập source mà principal được đọc rồi filter (\texttt{principal} $\rightarrow$ \texttt{sources}).
\end{itemize}

Một principal không bắt buộc là một người - có thể là một người, một service, hoặc một định danh theo nhóm chức năng; độ mịn của định danh do cách cấp phát credential quyết định (một credential cho mỗi vai trò, hoặc một credential cho mỗi người).

\begin{itemize}
  \item \textbf{FR-8.1} Hệ thống phải phục vụ nhiều AI Agent truy vấn đồng thời trên cùng một kho tri thức dùng chung và trả kết quả nhất quán.
  \item \textbf{FR-8.2} Hệ thống phải xác thực (authenticate) các request tới REST API và MCP, trước khi cho phép truy cập tri thức.
  \item \textbf{FR-8.3} Hệ thống phải kiểm soát truy cập tri thức ở mức source thông qua một ACL do admin cấu hình, xác định principal nào được đọc (\texttt{read}) tri thức từ source nào. Đơn vị phân quyền là source (\texttt{source\_id}).
      
  \item \textbf{FR-8.4} ACL phải gồm các thành phần:
  \begin{itemize}
  \item \textbf{Principal registry} -- mỗi principal gồm: \texttt{principal\_id}, \texttt{type} (\texttt{subject} | \texttt{group}; \texttt{subject} là một người, một service, hoặc một định danh dùng chung theo nhóm chức năng như \texttt{backend}), \texttt{members} (danh sách \texttt{principal\_id}, chỉ với \texttt{group}), \texttt{credentials} (danh sách \texttt{credential\_id} gắn với principal, không lưu giá trị secret thô) và \texttt{role} (\texttt{admin} | \texttt{member}).
  \item \textbf{Grants} -- mỗi grant gồm: \texttt{principal\_id}, danh sách \texttt{source\_id} áp dụng, và \texttt{permission} = \texttt{read}.
  \item \textbf{default\_policy} $=$ \texttt{deny} | \texttt{allow}.
  \end{itemize}
  \item \textbf{FR-8.5} Trên mỗi request, hệ thống phải resolve credential đã xác thực về đúng một principal. Quyền hiệu lực của một principal là hợp (union) quyền của chính nó và mọi \texttt{group} chứa nó (đệ quy).
  \item \textbf{FR-8.6} Một phản hồi truy xuất chỉ được chứa tri thức từ những source mà principal được đọc. Không một kết quả, node hay quan hệ kề trong đồ thị, hay metadata nào bắt nguồn từ source ngoài quyền đọc được xuất hiện trong phản hồi, kể cả khi nó tới gián tiếp qua duyệt đồ thị hoặc chỉ nổi lên sau khi hợp nhất.
  \item \textbf{FR-8.7} Quy tắc quyết định quyền đọc một source: dưới \texttt{default\_policy = deny}, principal chỉ đọc source được grant tường minh; dưới \texttt{default\_policy = allow}, principal đọc mọi source trừ source xuất hiện trong bất kỳ grant nào (source đó trở thành restricted, chỉ principal được grant mới đọc được).
  \item \textbf{FR-8.8} Hệ thống phải fail-closed: khi không resolve được principal, hoặc khi không xác định được quyền, request bị từ chối.
  \item \textbf{FR-8.9} Admin có thể tạo/sửa/xóa principal, group, grant và \texttt{default\_policy} qua REST API.
  \item \textbf{FR-8.10} Admin phải phát hành (issue) và thu hồi (revoke) credential gắn với một principal qua REST API; một principal có thể có nhiều credential.
  \item \textbf{FR-8.11} Giá trị secret của credential chỉ hiển thị một lần lúc phát hành và lưu ở dạng không thể khôi phục.
  \item \textbf{FR-8.12} Credential bị thu hồi phải bị từ chối ngay từ request kế tiếp.
  \item \textbf{FR-8.13} Các thao tác quản trị -- quản lý principal/group/grant/\texttt{default\_policy} (FR-8.9), cấp/thu hồi credential (FR-8.10), và cấu hình nguồn dữ liệu (FR-1.5, FR-5.1) -- chỉ principal có \texttt{role = admin} mới được thực hiện; request không đủ quyền bị từ chối.
\end{itemize}

\subsection{Yêu cầu phi chức năng (Non-functional requirements)}

\subsubsection{NFR-1: Hiệu năng truy xuất (Performance / Latency)}

\begin{itemize}
  \item \textbf{NFR-1.1} Với truy vấn top-k (k $\leq$ 10) trên kho đã index, độ trễ p95 phải $\leq$ 2 giây và p50 $\leq$ 0.8 giây, không tính độ trễ của embedding provider bên ngoài.
  \item \textbf{NFR-1.2} Với truy vấn trùng nội dung đã có trong cache, hệ thống phải trả kết quả trong $\leq$ 100 ms.
\end{itemize}

\subsubsection{NFR-2: Chất lượng truy xuất (Retrieval quality)}

\begin{itemize}
  \item \textbf{NFR-2.1} Trên một bộ truy vấn kiểm thử có nhãn (gold set $\geq$ 50 truy vấn), hệ thống phải đạt Recall@10 $\geq$ 0.85 và MRR $\geq$ 0.70.
  \item \textbf{NFR-2.2} Trên cùng bộ kiểm thử, hybrid search phải cải thiện Recall@10 ít nhất 5 điểm phần trăm so với chỉ dùng semantic search.
\end{itemize}

\subsubsection{NFR-3: Khả năng mở rộng (Scalability)}

\begin{itemize}
  \item \textbf{NFR-3.1} Khi mở rộng theo chiều ngang bằng cách bổ sung hoặc thay thế backend lưu trữ, hệ thống phải chạy đúng và vượt qua toàn bộ test chức năng.
\end{itemize}

\subsubsection{NFR-4: Hiệu quả cập nhật \& độ mới (Update efficiency / Freshness)}

\begin{itemize}
  \item \textbf{NFR-4.1} Cập nhật incremental cho một commit ($\leq$ 50 file thay đổi) phải hoàn tất trong $\leq$ 2 phút.
\end{itemize}

\subsubsection{NFR-5: Độ tin cậy \& ổn định (Reliability)}

\begin{itemize}
  \item \textbf{NFR-5.1} Khi nạp lỗi một file hoặc nguồn, hệ thống phải bỏ qua phần lỗi, giữ nguyên chỉ mục hiện có và ghi log; kiểm chứng bằng test inject lỗi.
  \item \textbf{NFR-5.2} Lệnh gọi phụ thuộc bên ngoài (embedding provider) phải retry tối thiểu 3 lần kèm backoff trước khi báo lỗi; mọi lỗi phải được log.
\end{itemize}

\subsubsection{NFR-6: Bảo mật \& quyền riêng tư (Security \& privacy)}

\begin{itemize}
  \item \textbf{NFR-6.1} 100\% endpoint REST và MCP phải yêu cầu xác thực (FR-8.2); request không có khóa hợp lệ phải bị từ chối.
  \item \textbf{NFR-6.2} Không được hard-code bất kỳ secret nào (API key, mật khẩu DB) trong mã nguồn; kiểm chứng bằng quét secret.
  \item \textbf{NFR-6.3} Với một principal không có quyền đọc source S, 0\% kết quả truy xuất -- kể cả kết quả đến từ graph traversal và sau fusion -- có nguồn gốc từ S; kiểm chứng bằng bộ test phủ từng đường truy xuất (semantic, keyword, graph).
  \item \textbf{NFR-6.4} Việc áp ACL không được làm tăng độ trễ p95 của truy vấn quá 10\% so với cùng truy vấn khi không áp ACL, trên cùng kho và cùng mức tải.
  \item \textbf{NFR-6.5} Mọi thay đổi config runtime (ACL, danh sách source, tunable) phải có hiệu lực với mọi request nhận sau thời điểm thay đổi được ghi nhận mà không cần restart, với độ trễ áp dụng $\leq$ 5 giây; kiểm chứng bằng test thay đổi rồi truy vấn lại.
\end{itemize}

\subsubsection{NFR-7: Tính dễ dùng của API (API usability)}

\begin{itemize}
  \item \textbf{NFR-7.1} 100\% endpoint phải có mô tả OpenAPI kèm ví dụ request/response (FR-5.3).
\end{itemize}

\subsubsection{NFR-8: Đóng gói \& triển khai (Deployability \& portability)}

\begin{itemize}
  \item \textbf{NFR-8.1} Mọi thành phần của hệ thống phải có image Docker.
  \item \textbf{NFR-8.2} Toàn hệ thống phải khởi chạy bằng một lệnh \texttt{docker compose up} và sẵn sàng nhận request trong $\leq$ 5 phút trên máy chuẩn.
  \item \textbf{NFR-8.3} Chuyển giữa các môi trường (dev, server, cloud) chỉ bằng thay đổi cấu hình, không build lại image; kiểm chứng bằng việc chạy đúng ở cả ba.
\end{itemize}

\subsubsection{NFR-9: Khả năng bảo trì (Maintainability)}

\begin{itemize}
  \item \textbf{NFR-9.1} Thay embedding provider chỉ bằng thay đổi cấu hình (FR-7.1); tiêu chí: chuyển đổi thành công và vượt qua bộ test chức năng.
  \item \textbf{NFR-9.2} Hệ thống phải log mọi sự kiện ingestion, truy vấn và lỗi; tiêu chí: tái dựng lại được một request lỗi chỉ từ log.
  \item \textbf{NFR-9.3} Thêm một định dạng tài liệu mới hoặc một nguồn dữ liệu mới chỉ cần bổ sung một thành phần xử lý độc lập; kiểm chứng bằng việc thêm một parser mẫu.
\end{itemize}



